\begin{center}
    \begin{singlespace}
        \includegraphics[width=2cm]{logo_USB}

        {\large
          UNIVERSIDAD SIMÓN BOLÍVAR\\
        }

        \textbf{
          DECANATO DE ESTUDIOS PROFESIONALES\\
          COORDINACIÓN DE INGENIERÍA DE LA COMPUTACIÓN
        }

        \vspace{0.5cm}
        \textbf{
          DESARROLLO DEL MÓDULO PRINCIPAL Y ESTADÍSTICAS DE LA LIBRERÍA AUDITORÍAS TURPIAL
        }

        \vspace{0.5cm}
        INFORME DE PASANTÍA LARGA\\
        Realizado por:
        Stefani Carolina Castellanos Torres\\

        Con la asesoría de: Prof. Angela Di Serio, Ing. Pedro Romero.\\
    \end{singlespace}
\end{center}

Auditorías Turpial es una librería para Django que permite mantener un historial de las operaciones realizadas sobre alguna tabla de base de datos, a través de auditorías. La empresa Turpial Development, requiere una librería de Django que le permita acelerar el desarrollo de aplicaciones que necesiten un sistema de este tipo, y esté desarrollado bajo la arquitectura de microservicios para que sea fácil de integrar. Esta librería consta de tres módulos: Principal (\textit{Core}), Estadísticas y Reportes. Adicionalmente, estos microservicios pueden ser actualizados mediante una herramienta de Integración Continua. En el presente informe, se describen los procesos de diseño, implementación y pruebas de parte del módulo \textit{Core} y completamente el módulo Estadísticas de Auditorías Turpial, desarrollado como pasantía larga en Turpial Development siguiendo la metodología \textit{Turpial Agile Unified Process}. Así como, la implantación del sistema de Integración Continua utilizando Jenkins. Como resultado, se lograron todos los objetivos planteados y otros originados durante la fase de construcción. \\

\textbf{Palabras claves}: Auditoría de sistemas, librería, Django, Integración Continua, microservicio.
